\documentclass[10pt,a4paper]{amsart}
\usepackage[margin=19mm]{geometry}
\usepackage{amsmath,amssymb,mathtools}
\usepackage[scr=rsfs,cal=euler]{mathalfa}
\usepackage{xfrac}
\usepackage[unicode,hidelinks,bookmarks=false]{hyperref}
\newcommand{\ie}{i.e.\ }
\newcommand{\cf}{cf.\ }
\newcommand{\resp}{respectively\ }
\newcommand{\Subsection}{\subsection}
\providecommand{\nobreakdash}{\nobreak\hbox{-}}
\setlength{\emergencystretch}{2em}
\multlinegap0pt
\usepackage{bm}
\usepackage{bbm}
\usepackage{dsfont}
\usepackage{xspace}
\usepackage{cancel}
\usepackage{mathtools}
\usepackage{amsthm}
\makeatletter
\@ifundefined{Theorem}{\newtheorem{Theorem}{Theorem}}{}
\@ifundefined{Lemma}{\newtheorem{Lemma}[Theorem]{Lemma}}{}
\@ifundefined{Proposition}{\newtheorem{Proposition}[Theorem]{Proposition}}{}
\makeatother
\title[Characterization of Gaussian Tensor Ensembles]{Characterization of Gaussian Tensor Ensembles}
\author{R\'emi Bonnin}
\date{Source: arXiv:2505.02814v2 (CC BY 4.0)}
\begin{document}
\maketitle

\noindent\emph{Source and licence.} Characterization of Gaussian Tensor Ensembles. Version arXiv:2505.02814v2 (CC BY 4.0), distributed under the Creative Commons Attribution 4.0 International licence, as recorded in the arXiv licence metadata for this exact version and shown on the abstract page https://arxiv.org/abs/2505.02814v2. Full rights notice: licenses/arxiv-2505.02814-CC-BY-4.0.txt.\par\smallskip

\noindent\textit{Editorial scope.} Counted unit: the Lemma whose source label is \texttt{xy} in the article (source0.tex lines 344--350, source label \texttt{xy}), delivered together with the article's own complete original proof of it (source0.tex lines 352--392, one \texttt{proof} environment of 41 lines). No mathematics is written, repaired or re-derived: statement and proof are the article's own text line for line. Required context retained uncounted: maxwell (lines 289--296); derivatives (lines 337--342). Every definition, display and label of those lines is retained unchanged. Omitted from this package and not redistributed: 1--288, 297--336, 343--343, 351--351, 393--873. Nothing inside a retained excerpt is omitted or altered: every retained body is delivered exactly as the article's lines are written, so no omission receipt of any kind accompanies this package. Citations: the retained excerpts carry no citation command at all, so no bibliography entry of the article is transcribed and no reference is redistributed. Reference adaptation: none was needed and none was made; every reference inside the delivered unit resolves inside it, so no unresolved reference or citation is produced. Printed slips kept literal: no printed word, symbol, hypothesis or number of the counted statement or of its proof was altered. Layout and class: the article's own class is \texttt{sigma} with packages including ifpdf, dsfont, tikz; the wrapper is the lane's pinned \texttt{amsart} editorial wrapper at 10pt on A4, which restates the theorem-like environments the retained text needs and adds the article's own macro definitions that the retained text needs (none), with extra wrapper packages (bm, bbm, dsfont, xspace, cancel, mathtools). The delivered numbering is the unit's own. Assets: no retained span contains a figure environment, a graphics-inclusion command, a PDF-page inclusion, a table or a TikZ picture; no image file is copied into this package, the package declares no asset dependency and no image file is redistributed with it. Licence: version arXiv:2505.02814v2 of this article is distributed under the Creative Commons Attribution 4.0 International licence, as recorded in the arXiv licence metadata for this exact version and shown on the abstract page https://arxiv.org/abs/2505.02814v2. A full rights notice is delivered at licenses/arxiv-2505.02814-CC-BY-4.0.txt. Characters: every retained excerpt is pure ASCII, so the delivered engine, XeLaTeX with its default Latin Modern text font, reports no missing character.
\begin{Theorem}[main]\label{maxwell}
For any $p\geq 1$ and $N\geq 2$, a real symmetric $($resp.\ hermitian, resp.\ self-dual hermitian$)$ random tensor $H$ of order $p$ and dimension $N$ has in the same time independent entries $($up to symmetries$)$ and a law invariant by orthogonal $($resp.\ unitary, resp.\ symplectic$)$ transformation if and only if it belongs to the Gaussian orthogonal $($resp.\ unitary, resp.\ symplectic$)$ tensor ensemble, that is, it has a law with a density of the form
\[ H \mapsto \alpha \exp \left( - \| H - \beta \mathcal{I} \|_{\mathrm{F}}^2 / \gamma \right),\]
where $(\alpha,\beta,\gamma) \in \mathbb{R}^*_+ \times \mathbb{R} \times \mathbb{R}^*_+$,
or equivalently,
\[ H \mapsto \exp \left( - a \| H \|_{\mathrm{F}}^2 + b \operatorname{Tr}^{\bullet \bullet}(H) + c \right),\]
where $(a,b,c) \in \mathbb{R}^*_+ \times \mathbb{R} \times \mathbb{R}$.
\end{Theorem}

\begin{Proposition}\label{derivatives}
We have
    \[ \left( \frac{\partial H}{\partial \theta} \right) = \sum_{r=1}^p A \star_1 H^{(1,r)}. \]
In other words,
\[\left( \frac{\partial H}{\partial \theta} \right)_{i_1,\dots,i_p} = \sum_{k=1}^N (A_{i_1k}H_{k,i_2,\dots,i_p} + A_{i_2k}H_{i_1,k,\dots,i_p} + \dots + A_{i_pk}H_{i_1,\dots,i_{p-1},k}).\]
\end{Proposition}

\begin{Lemma}\label{xy}
Let $g_1$, $g_2$, $g_3$ be three continuous and differentiable functions such that
$ g_1(xy)=g_2(x)+g_3(y)$.
Then, there exists $\alpha$, $\beta_1$, $\beta_2$, $\beta_3$ such that for all $k\in \{1,2,3\}$,
$g_k(u)=\alpha \ln u + \beta_k$,
with $\beta_1 = \beta_2+\beta_3$.
\end{Lemma}

\begin{proof}[Proof of Theorem~\ref{maxwell}.]
Let $H$ be a real symmetric tensor having independent entries and a~law $\mu$ orthogonally invariant. Possibly after changing $H$ into $H+\epsilon W$ with $W$ belonging to the GOTE, we may assume without loss of generality that $\mu$ has a smooth density given by
\[ f(H)=\prod_{i_1\leq \dots \leq i_p} f_{i_1,\dots,i_p}(H_{i_1,\dots,i_p}). \]
By orthogonal invariance,
$\frac{\partial}{\partial \theta} \ln f(H)=0$.
Then using the computation of Proposition \ref{derivatives}, we find
\begin{align*}
    0 &= \sum_{i_1\leq \dots \leq i_p} \frac{1}{f_{i_1,\dots,i_p}(H_{i_1,\dots,i_p})} \frac{\partial f_{i_1,\dots,i_p}}{\partial H_{i_1,\dots,i_p}} \frac{\partial H_{i_1,\dots,i_p}}{\partial \theta} \\
    &= \sum_{i_1\leq \dots \leq i_p} \frac{1}{f_{i_1,\dots,i_p}} \frac{\partial f_{i_1,\dots,i_p}}{\partial H_{i_1,\dots,i_p}} \sum_{k=1}^N (A_{i_1k}H_{k,i_2,\dots,i_p} + A_{i_2k}H_{i_1,k,\dots,i_p} + \dots + A_{i_pk}H_{i_1,\dots,i_{p-1},k}).
\end{align*}
Moreover, we know that $A_{i,k}$ is equal to $-1$ if $(i,k)=(1,2)$, to $1$ if $(i,k)=(2,1)$ and to $0$ otherwise. We introduce a brief notation for ease of notation in the sums with multi-indices. For given integers $r \in \{1,\dots, p-1\}$ and $i_{r+1},\dots,i_p$, let us denote $f_{a_1,\dots,a_r}$ for $f_{a_1,\dots,a_r,i_{p-r},\dots,i_p}$ and $H_{a_1,\dots,a_r}$ for $H_{a_1,\dots,a_r,i_{p-r},\dots,i_p}$. In particular, $f_1(H_1)=f_{1,i_{2},\dots,i_p}(H_{1,i_{2},\dots,i_p})$. The previous equations gives finally
\begin{align}
    0= {}&\sum_{3\leq i_2\leq \dots \leq i_p} \left[ \frac{-1}{f_1(H_1)} \frac{\partial f_{1}}{\partial H_{1}} H_2 + \frac{1}{f_2(H_2)} \frac{\partial f_{2}}{\partial H_{2}} H_1 \right]\nonumber\\
    &+ \sum_{3\leq i_3\leq \dots \leq i_p} \left[ 2 \left( \frac{-1}{f_{1,1}(H_{1,1})} \frac{\partial f_{1,1}}{\partial H_{1,1}} + \frac{1}{f_{2,2}(H_{2,2})} \frac{\partial f_{2,2}}{\partial H_{2,2}} \right) H_{1,2}\right.\nonumber\\
    & \left.+ \frac{1}{f_{1,2}(H_{1,2})} \frac{\partial f_{1,2}}{\partial H_{1,2}} (H_{1,1} - H_{2,2}) \right]+ \sum_{3\leq i_4\leq \dots \leq i_p} \left[ \dots \right].\label{eq-A}
\end{align}
The different brackets depend on mutually exclusive sets of variables, and their sum is $0$, so each must be a constant and in particular for given integers $i_p\geq \dots \geq i_2 \geq 3$, there exists~${C=C_{i_2,\dots,i_p}}$ such that
\[C = \frac{-1}{f_1(H_1)} \frac{\partial f_{1}}{\partial H_{1}} H_2 + \frac{1}{f_2(H_2)} \frac{\partial f_{2}}{\partial H_{2}} H_1. \]
Hence, after dividing by $H_1H_2$ we get
\[ \frac{C}{H_1H_2} = \frac{-1}{H_1} \frac{1}{f_1(H_1)} \frac{\partial f_{1}}{\partial H_{1}} + \frac{1}{H_2}\frac{1}{f_2(H_2)} \frac{\partial f_{2}}{\partial H_{2}}, \]
and now applying Lemma \ref{xy} we find that the constant $C$ must be zero. That means
\[ \frac{1}{H_1} \frac{1}{f_1(H_1)} \frac{\partial f_{1}}{\partial H_{1}} = \frac{1}{H_2}\frac{1}{f_2(H_2)} \frac{\partial f_{2}}{\partial H_{2}}, \]
and this must be a constant $c'=c'_{i_2,\dots,i_p}$, which gives
\[\frac{1}{f_1(H_1)} \frac{\partial f_{1}}{\partial H_{1}} = c' H_1.\]
After integrating we find $\ln f_1(H_1) = \frac{c'}{2} H_1^2$, so $c'$ must be negative ($f_1$ is a density) and we denote it $-2a$ in place of, with $a \in \mathbb{R}^*_+$ (depending on ${i_2,\dots,i_p}$). Finally, we conclude from all of this
\smash{$f_{1,i_{2},\dots,i_p}(H_{1,i_{2},\dots,i_p}) = \exp\bigl(-a H_{1,i_{2},\dots,i_p}^2\bigr)$}.
The same holds for $f_2(H_2)$ giving $f_{2,i_{2},\dots,i_p}(H_{2,i_{2},\dots,i_p}) = \exp\bigl(-a H_{2,i_{2},\dots,i_p}^2\bigr)$.
Note that involving the symmetries and orthogonal matrices $U_{\theta}$ with the rotation acting at different positions than~$(1,2)$, it follows that any $p$-tuple of indices with at least {\em one index $i_k$ distinct from all other ones} gives a contribution~${f_{i_{1},\dots,i_p}(H_{i_{1},\dots,i_p}) = \exp\bigl(-a_{i_1,\dots,i_{k-1},i_{k+1},\dots,i_p} H_{i_{1},\dots,i_p}^2\bigr)}$. Now, the second bracket in equation~\eqref{eq-A} gives for $i_p\geq\dots\geq i_3\geq 3$,
\[ 2 \left( \frac{-1}{f_{1,1}(H_{1,1})} \frac{\partial f_{1,1}}{\partial H_{1,1}} + \frac{1}{f_{2,2}(H_{2,2})} \frac{\partial f_{2,2}}{\partial H_{2,2}} \right) -2a_{1,2,i_3,\dots,i_p} (H_{1,1} - H_{2,2}) = \frac{C_{i_3,\dots,i_p}}{H_{1,2}} =0,\]
and we deduce immediately
\[ \frac{1}{f_{1,1}(H_{1,1})} \frac{\partial f_{1,1}}{\partial H_{1,1}} + a H_{1,1} = \frac{1}{f_{2,2}(H_{2,2})} \frac{\partial f_{2,2}}{\partial H_{2,2}} + a H_{2,2} = b_{i_3,\dots,i_p},\]
that is, finally
\[ f_{1,1,i_{3},\dots,i_p}(H_{1,1,i_{3},\dots,i_p}) = \exp\left(-\frac{a}{2} H_{1,1,i_{3},\dots,i_p}^2-bH_{1,1,i_{3},\dots,i_p}\right). \]
But as all invariants are expressible in terms of the trace invariants, the two only ones which can appear are the melon (corresponding to the Frobenius norm) and the bouquet (corresponding to the paired trace), and the constants must satisfy \smash{$a_{i_1,\dots,i_p}=\frac{a_{1,\dots,p}}{\Gamma_{i_1,\dots,i_p} }$}, $b_{i_1,i_1,\dots,i_{p/2},i_{p/2}}=\smash{\frac{b}{\Gamma_{i_1,i_1,\dots,i_{p/2},i_{p/2}} }}$ and $b_{i_1,\dots,i_p}=0$ otherwise, leading to
\begin{align*}
    f(H)={}& \exp \Bigg( -\sum_{i_1\leq\dots\leq i_p} \frac{a}{\Gamma_{i_1,\dots,i_p} } H_{i_{1},\dots,i_p}^2 \\
    &+ \mathds{1}_{p \text{ even}} \sum_{i_1 \leq\dots \leq i_{p/2}} \frac{b}{\Gamma_{i_1,i_1,\dots,i_{p/2},i_{p/2}} } H_{i_1,i_1,\dots,i_{p/2},i_{p/2}} +c \Bigg) \\
    ={}& \exp \Bigg( -a \sum_{i_1,\dots,i_p} H_{i_{1},\dots,i_p}^2 + b \mathds{1}_{p \text{ even}} \sum_{i_1,\dots, i_{p/2}} H_{i_1,i_1,\dots,i_{p/2},i_{p/2}} +c \Bigg).
\end{align*}
This concludes the proof.
\end{proof}

\end{document}
